\documentclass[9pt,a4paper,twocolumn,twoside]{rho-class/rho}

% --- 1. Idioma y matemáticas básicas ---
\usepackage[spanish,es-nodecimaldot,es-noindentfirst,es-noshorthands]{babel}


% --- 2. EL TRUCO DE MAGIA (BIEN ESCRITO) ---
% Borramos los comandos problemáticos usando su nombre exacto interno
\expandafter\let\csname not=\endcsname\relax
\expandafter\let\csname not<\endcsname\relax
\expandafter\let\csname not>\endcsname\relax

% Ahora ya podemos cargar la fuente moderna sin que explote
\usepackage{fontspec}
\usepackage{unicode-math}

% --- 3. Bibliografía ---
\addbibresource{rho.bib}

% --- 4. Datos del Informe ---
\title{Estudio de la radiación con el experimento de Geiger-Müller}
\author{Edgar Fernández Vigil}
\affil{Física - Facultad de Ciencias - Universidad de Oviedo}
\dates{\today}
\journalname{TEIII - Práctica de Laboratorio}
\journal{Radiación con el experimento de Geiger-Müller}

% --- 5. Resumen ---
\begin{abstract}
En el presente informe se detalla el estudio experimental de las propiedades de la radiación ionizante y su interacción con la materia mediante el uso de un tubo contador Geiger-Müller. %
En primer lugar, se caracterizó la respuesta instrumental determinando empíricamente su curva \textit{plateau} para establecer el voltaje óptimo de operación. %
Posteriormente, se analizó la naturaleza estocástica de las emisiones radiactivas, verificando que la detección de la radiación de fondo en intervalos cortos obedece a una distribución asimétrica de Poisson, mientras que las tasas de conteo elevadas procedentes de fuentes activas se ajustan a una distribución continua de Gauss. %
Asimismo, se comprobó la atenuación geométrica de una fuente isótropa de $^{90}\text{Sr}$, confirmando su decaimiento proporcional a la inversa del cuadrado de la distancia. %
Por último, se evaluó la atenuación material de la radiación; mediante la interposición de láminas de aluminio y plomo, se calculó el coeficiente de absorción y el rango de frenado máximo para partículas $\beta$, además de contrastar cualitativamente el poder de penetración y las necesidades de apantallamiento para emisiones $\alpha$ ($^{238}\text{Pu}$) y $\gamma$ ($^{60}\text{Co}$). %
\end{abstract}

% ----------------------------------------------------------
\begin{document}

    \maketitle
    \thispagestyle{firststyle}

\section{Introducción}

\rhostart{E}l tubo Geiger-Müller es uno de los instrumentos fundamentales en la física experimental para la detección de radiación ionizante en entornos de baja actividad. %
Su funcionamiento se basa en la ionización de un gas encerrado en un cilindro sometido a una alta diferencia de potencial; cuando una partícula o fotón incidente penetra en la región activa, desencadena una avalancha de electrones que se registra como un pulso eléctrico. %
Para operar de manera fiable este dispositivo, es imperativo caracterizar su respuesta frente al voltaje aplicado, determinando su curva \textit{plateau} para establecer un voltaje óptimo de trabajo y evitar la región de descarga continua. %

Dado que la desintegración radiactiva es un proceso fundamentalmente estocástico, la detección de eventos en intervalos de tiempo fijos obedece a leyes estadísticas precisas. %
En regímenes de baja tasa de conteo, como ocurre al medir la radiación de fondo ambiental (compuesta, entre otros, por rayos cósmicos y muones), el proceso obedece a una distribución asimétrica de Poisson. %
Por el contrario, cuando se expone el detector a fuentes radiactivas con tasas de emisión elevadas, el comportamiento límite permite aproximar las fluctuaciones estadísticas a una distribución continua y simétrica de Gauss. %

Adicionalmente, el estudio de la radiación exige comprender su propagación espacial y su interacción con la materia. %
Por un lado, la fluencia de partículas emitidas por una fuente puntual isótropa se atenúa geométricamente con el cuadrado de la distancia. %
Por otro lado, la atenuación material depende drásticamente de la naturaleza de la radiación incidente: las partículas $\beta$ (electrones) sufren un frenado paulatino caracterizado por un rango máximo y un coeficiente de absorción de masa, mientras que la radiación $\alpha$ y $\gamma$ presentan mecanismos de interacción y requisitos de apantallamiento muy diferentes al atravesar materiales como el aluminio o el plomo. %

\section{Objetivos}

La finalidad general de esta práctica es familiarizarse con los métodos de medida de radiación en baja intensidad y comprender los efectos de la interacción radiación-materia. %
De forma específica, el experimento persigue los siguientes objetivos: %

\begin{itemize}
    \item \textbf{Caracterización del tubo Geiger-Müller:} Obtener experimentalmente la curva \textit{plateau} del detector registrando la tasa de conteo en función del voltaje aplicado, con el fin de determinar el punto de operación óptimo (en torno a 1000 V) y calcular la pendiente de dicha meseta. %
    \item \textbf{Análisis estadístico de la emisión:} Contrastar la naturaleza estocástica de la radiactividad verificando empíricamente que la radiación de fondo obedece a una distribución de Poisson, mientras que las medidas de una muestra radiactiva activa (como el $^{90}\text{Sr}$ o $^{137}\text{Cs}$) se ajustan a una campana de Gauss. %
    \item \textbf{Verificación de la atenuación geométrica:} Demostrar que la tasa de detección de una fuente radiactiva isótropa disminuye en proporción inversa al cuadrado de la distancia fuente-detector ($\sim 1/d^2$). %
    \item \textbf{Estudio de la absorción de partículas $\beta$:} Medir la atenuación continua de los electrones emitidos por una fuente de $^{90}\text{Sr}$ al interponer láminas de aluminio y plomo de distinto grosor, para calcular su rango máximo ($R$), su coeficiente de absorción y estimar la energía máxima teórica de la desintegración. %
    \item \textbf{Apantallamiento y detección de radiación $\alpha$ y $\gamma$:} Analizar el poder de penetración y la eficiencia relativa del detector frente a emisores $\alpha$ ($^{238}\text{Pu}$) y fotones $\gamma$ ($^{60}\text{Co}$ o $^{133}\text{Ba}$), comprobando la eficacia atenuadora de materiales como el papel, el aluminio y el plomo. %
\end{itemize}

\section{Procedimiento Experimental}

La toma de datos se ha estructurado en varias fases para estudiar tanto la respuesta instrumental del detector como las propiedades intrínsecas de las diferentes radiaciones. % [cite: 509]

El desarrollo de la práctica se dividió en las siguientes etapas principales:

\subsection{Determinación de la curva \textit{plateau}}
Para establecer el voltaje óptimo de operación del tubo Geiger-Müller, se midió la tasa de conteo de diversas muestras ($^{90}\text{Sr}$, $^{204}\text{Tl}$, $^{60}\text{Co}$ y $^{238}\text{Pu}$) variando la tensión de alimentación. % [cite: 525, 535]
Partiendo de 700 V (voltaje de arranque), se incrementó el potencial en pasos de 25 V hasta un máximo de 1200 V, registrando los pulsos en intervalos de 30 segundos. % [cite: 526, 533]
A partir de la representación gráfica de estos datos, se identificó la región de \textit{plateau} y se calculó su pendiente (que debe ser inferior al 10\% de incremento por cada 100 V), fijando finalmente el voltaje de trabajo en torno a los 1000 V. % [cite: 527, 528]

\subsection{Análisis estadístico (Distribuciones de Poisson y Gauss)}
Con el voltaje de operación fijado, se estudió el carácter estocástico de la radiación bajo dos regímenes distintos:
\begin{itemize}
    \item \textbf{Radiación de fondo (Poisson):} Se tomaron más de 100 medidas de corta duración (2 s) en ausencia de muestra radiactiva. % [cite: 577, 588]
    Al tratarse de una tasa de conteo muy baja (procedente principalmente de muones y radiación cósmica), los datos se agruparon en un histograma de frecuencias y se ajustaron a una distribución discreta y asimétrica de Poisson, verificando que la varianza coincidiera estadísticamente con la media. % [cite: 560, 577, 581]
    \item \textbf{Radiación de la muestra (Gauss):} Se introdujo una fuente activa ($^{90}\text{Sr}$ o $^{137}\text{Cs}$) y se registraron más de 500 medidas de 1 segundo de duración. % [cite: 627, 628, 631]
    Dado el elevado número de sucesos independientes, el histograma resultante se ajustó a una distribución continua y simétrica de Gauss. % [cite: 620, 621, 623]
\end{itemize}

\subsection{Dependencia geométrica con la distancia}
Para comprobar la ley de la inversa del cuadrado, se utilizó la fuente de $^{90}\text{Sr}$ (asumida como puntual e isótropa) y se varió sistemáticamente la distancia entre esta y la ventana del detector en pasos de 1 cm. % [cite: 633, 634, 640]
Se tomaron 5 medidas de 10 segundos en cada escalón espacial. % [cite: 639]
Los datos de actividad frente a distancia se ajustaron analíticamente considerando el ángulo sólido del detector, lo que permitió además estimar la actividad inicial de la muestra y, por consiguiente, su antigüedad. % [cite: 635, 636, 637]

\subsection{Absorción y apantallamiento en la materia}
Finalmente, se evaluó el poder de penetración de las distintas radiaciones al atravesar materiales absorbentes:
\begin{itemize}
    \item \textbf{Partículas $\beta$:} Se colocó la fuente de $^{90}\text{Sr}$ en una posición fija y se interpusieron láminas de aluminio y plomo de densidad superficial conocida. % [cite: 654, 656]
    Tomando varias medidas por lámina, se obtuvo una curva de transmisión exponencial que permitió deducir el coeficiente de absorción y el rango máximo ($R$) en la aproximación de frenado continuo (CSDA). % [cite: 653, 655]
    A partir de este rango, se estimó la energía máxima teórica de las partículas $\beta$ emitidas. % [cite: 659]
    \item \textbf{Radiación $\alpha$ y $\gamma$:} Se comprobó la alta capacidad de penetración de los fotones $\gamma$ frente a la exigua capacidad de las partículas $\alpha$, midiendo la tasa de atenuación al interponer barreras de papel, aluminio y plomo ante fuentes puras de $^{238}\text{Pu}$ y $^{60}\text{Co}$. % [cite: 666, 668]
\end{itemize}


\section{Resultados y Discusión}
holaaa

\end{document}
